\documentclass[11pt]{article}
\usepackage[margin=0.9in]{geometry}
\usepackage{graphicx}
\usepackage{caption}
\usepackage{hyperref}
\usepackage{float}
\usepackage{parskip}
\usepackage{booktabs}
\usepackage{array}

\hypersetup{colorlinks=false, hidelinks}

% Compact layout: tighter paragraph gaps, small captions, denser headings.
\setlength{\parskip}{3pt plus 1pt}
\captionsetup{font=small, skip=3pt}
\captionsetup{labelfont=bf}

% Guarded diagram include: renders the PNG if present, else a labelled
% placeholder, so the document always compiles.
\newcommand{\dgram}[4]{%
  \begin{figure}[htbp]
    \centering
    \IfFileExists{figures/#1}{%
      \includegraphics[width=#2\textwidth,height=0.33\textheight,keepaspectratio]{figures/#1}%
    }{%
      \fbox{\parbox[c][3cm][c]{0.7\textwidth}{\centering\itshape
        [\,Diagram placeholder\,]\\[6pt]Add \texttt{report/figures/#1}.}}%
    }
    \caption{#3}\label{#4}
  \end{figure}%
}

\title{\textbf{GreenStep Employee Sustainability Challenge}\\[0.4em]
\large Technical Design Document}
\author{Oli Gurmessa \and Hadi Shaar \and Naisha Srivastav \and Nathanael Agbemadon}
\date{CISC480 --- Senior Capstone\\[0.3em]
Instructor: Dr.\ Jimenez \quad\textbullet\quad Spring 2026\\[0.3em]
Submitted: May 22, 2026}

\begin{document}

\maketitle
\tableofcontents
\newpage

% ─────────────────────────────────────────────────────────────────────────────
\section{System Overview / Proposed Solution}
% ─────────────────────────────────────────────────────────────────────────────

GreenStep Employee is a web-based, gamified sustainability tracking platform
designed for workplace use. Employees log sustainable actions across five
categories---Transport, Water, Energy, Recycling, and Food---earning points,
streaks, and badges that contribute to team and individual leaderboards. The
platform replaces manual Excel-based tracking with a real-time, low-friction
digital system.

The solution is delivered as a single web application accessible from any
browser. It requires no installation and is hosted on Vercel's free tier,
keeping it viable within the project's zero-budget constraint. Participation is
organized into \emph{challenges}: every user is auto-enrolled in a permanent
global challenge and may join additional, code-gated challenges.

% ─────────────────────────────────────────────────────────────────────────────
\section{Architecture Pattern}
% ─────────────────────────────────────────────────────────────────────────────

\subsection{Pattern: Modular Monolith (Layered)}

The system follows a \textbf{Modular Monolith} pattern organized into three
logical layers---Presentation, Application, and Data---all deployed as a single
Next.js~15 application (Figure~\ref{fig:arch}). The \textbf{Application layer}
is not a separate server: it comprises the Next.js \emph{API route handlers},
\emph{server actions} (\texttt{app/actions/challenge.ts}), and the
\emph{Clerk middleware} that authenticates every non-public request before it
reaches a route.

\dgram{arch_pattern.png}{0.95}{Modular-monolith layered architecture.
Presentation (React pages) calls the Application layer (API routes, server
actions, Clerk middleware), which calls the Data layer (Prisma ORM over Neon
PostgreSQL). Clerk is a cross-cutting external dependency that validates every
request.}{fig:arch}

\subsection{Design Decisions and Justifications}

Table~\ref{tab:just} links each major decision to a project requirement or
constraint.

\begin{table}[H]
\centering
\caption{Design decisions and their justifications.}
\label{tab:just}
\begin{tabular}{>{\bfseries}p{3.6cm} p{5.2cm} p{4.6cm}}
\toprule
Decision & Justification & Requirement/Constraint \\
\midrule
Modular monolith over microservices &
Single deployment unit; no inter-service networking or orchestration. &
4-person team, one semester, no DevOps budget. \\
\addlinespace
Next.js 15 for UI \emph{and} API &
Co-locating frontend, API routes, and server actions removes a separate
backend, CORS, and a second pipeline. &
Simplicity; single Vercel deployment; free tier. \\
\addlinespace
Clerk for authentication &
Production-grade auth, sessions, and JWTs without building or securing a
custom system; admin via \texttt{publicMetadata}. &
Security requirement; no auth-infra budget. \\
\addlinespace
Prisma + Neon PostgreSQL &
Type-safe access over a relational 7-model schema; serverless Postgres with a
free tier and no maintenance. &
Relational data; zero budget; low maintenance. \\
\addlinespace
Vercel hosting &
Zero-config build and deploy on push to \texttt{main}. &
Free tier; CI/CD without extra tooling. \\
\addlinespace
Five action categories (enum) &
Mirrors the partner's existing sustainability programme domains. &
Partner requirement. \\
\bottomrule
\end{tabular}
\end{table}

% ─────────────────────────────────────────────────────────────────────────────
\section{Architectural Views}
% ─────────────────────────────────────────────────────────────────────────────

\subsection{Logical View}

Figure~\ref{fig:logical} maps the three layers to concrete modules: the
\texttt{app/} pages (Presentation); the \texttt{app/api/*} route handlers,
\texttt{app/actions/*} server actions, \texttt{middleware.ts}, and
\texttt{lib/*} helpers (Application); and \texttt{lib/db.ts} with
\texttt{prisma/schema.prisma} over Neon (Data).

\dgram{logical_view.png}{0.95}{Logical view: the modules that make up each
layer of the application.}{fig:logical}

\subsection{Deployment View}

Figure~\ref{fig:deploy} shows the runtime nodes. The browser talks only to the
Next.js app on Vercel; the app talks to Clerk (HTTPS) and Neon
(TCP/PostgreSQL via Prisma).

\dgram{deployment_view.png}{0.9}{Deployment view: Vercel (Next.js app), Neon
(PostgreSQL), and Clerk (auth) as separate nodes, with the required
environment variables.}{fig:deploy}

\subsection{Data-Flow View}

Figure~\ref{fig:dataflow} traces a single request---logging an action---from
the browser through authentication, challenge resolution, the serializable
transaction, and the streak update.

\dgram{dataflow_log_action.png}{0.78}{Data-flow view for ``log an action,''
from \texttt{POST /api/actions} to the response.}{fig:dataflow}

% ─────────────────────────────────────────────────────────────────────────────
\section{C4 Model Diagrams}
% ─────────────────────────────────────────────────────────────────────────────

\subsection{Level 1 --- Context}
\dgram{c4_context.png}{0.85}{C4 Context. The GreenStep system used by Employee
and Admin actors, delegating authentication to Clerk and persistence to Neon
PostgreSQL.}{fig:ctx}

Figure~\ref{fig:ctx} gives the highest-level view: employees log actions and
view leaderboards; admins create challenges, approve teams, and export data.

\subsection{Level 2 --- Container}
\dgram{c4_container.png}{0.85}{C4 Container. One Next.js container serves UI,
API routes, server actions, and middleware; it uses the Prisma client to reach
Neon and the Clerk SDK to validate sessions.}{fig:cont}

\subsection{Level 3 --- Component}
\dgram{c4_component.png}{0.95}{C4 Component (inside the Next.js container):
Clerk middleware; the UI pages; the thirteen API route handlers; the
\texttt{challenge.ts} server actions; the \texttt{lib} helpers
(\texttt{team-requests}, \texttt{global-challenge}, \texttt{authz}); and the
Prisma client, which is the single point of database access.}{fig:comp}

As shown in Figure~\ref{fig:comp}, every inbound request first passes through
the Clerk middleware. Authenticated requests reach a page or one of the API
route handlers; all database access is centralized through the Prisma client.

\subsection{Level 4 --- Code (Entity-Relationship)}
\dgram{er_model.png}{0.98}{C4 Code level shown as the entity-relationship model
of the seven Prisma entities and their keys and relations.}{fig:er}

Figure~\ref{fig:er} is the authoritative data model and is reused as the UML
class/data-model diagram in Section~\ref{sec:uml}.

% ─────────────────────────────────────────────────────────────────────────────
\section{UML Diagrams}\label{sec:uml}
% ─────────────────────────────────────────────────────────────────────────────

\subsection{Class / Data-Model Diagram}

The class/data-model view is the seven-entity diagram in
Figure~\ref{fig:er}. The entities are \texttt{User}, \texttt{Challenge},
\texttt{UserChallenge}, \texttt{Team}, \texttt{TeamMember}, \texttt{Action},
and \texttt{TeamRequest}. \texttt{UserChallenge} and \texttt{TeamMember} are
join tables implementing the many-to-many relations User$\leftrightarrow$Challenge
and User$\leftrightarrow$Team. Two enums constrain string fields:
\texttt{ActionCategory} and \texttt{TeamRequestStatus}.

\subsection{Sequence Diagrams}

\dgram{seq_log_action.png}{0.95}{Sequence: logging an action. The route
authenticates via Clerk, ensures the global challenge exists, runs a
serializable transaction (duplicate check then create), and updates the
streak.}{fig:seq1}

\dgram{seq_team_request.png}{0.95}{Sequence: team request and admin approval.
A user submits a \texttt{PENDING} request; an admin approves it, triggering a
\texttt{\$transaction} that creates the Team and TeamMember and marks the
request \texttt{APPROVED}.}{fig:seq2}

Figure~\ref{fig:seq1} and Figure~\ref{fig:seq2} detail the two most
safety-critical flows: action logging (with its duplicate guard) and the
admin-gated team-creation workflow.

\subsection{Component Diagram}

The component view of the system---the Next.js internals and the external
Clerk and Neon dependencies---is given by the C4 Level~3 diagram,
Figure~\ref{fig:comp}.

% ─────────────────────────────────────────────────────────────────────────────
\section{Data Design and Representation}
% ─────────────────────────────────────────────────────────────────────────────

The schema (\texttt{prisma/schema.prisma}) defines seven models. Primary keys
are \texttt{cuid} strings. Tables~\ref{tab:core}--\ref{tab:join} list the
fields; ``U'' marks a unique constraint and ``FK'' a foreign key.

\begin{table}[H]
\centering
\caption{Core entities: User, Challenge, and Action.}
\label{tab:core}
\begin{tabular}{p{2.6cm} p{4.0cm} p{5.6cm}}
\toprule
Model & Key fields & Notable fields / relations \\
\midrule
\texttt{User} &
\texttt{id} (PK), \texttt{clerkId} (U), \texttt{email} (U) &
\texttt{name}, \texttt{department?}, \texttt{streak} (default 0),
\texttt{lastActionDate?}, \texttt{onboarded} (default false); has many
\texttt{Action}, \texttt{UserChallenge}, \texttt{TeamMember}, \texttt{TeamRequest}. \\
\addlinespace
\texttt{Challenge} &
\texttt{id} (PK), \texttt{code} (U) &
\texttt{name}, \texttt{startDate?}, \texttt{endDate?}; has many
\texttt{UserChallenge}, \texttt{Team}, \texttt{Action}, \texttt{TeamRequest}. \\
\addlinespace
\texttt{Action} &
\texttt{id} (PK) &
\texttt{userId} (FK), \texttt{challengeId?} (FK), \texttt{category}
(\texttt{ActionCategory}), \texttt{actionType}, \texttt{points} (default 0),
\texttt{note?}, \texttt{imageUrl?} (Text), \texttt{date}. \\
\bottomrule
\end{tabular}
\end{table}

\begin{table}[H]
\centering
\caption{Team, TeamRequest, and the two join tables.}
\label{tab:join}
\begin{tabular}{p{2.8cm} p{4.0cm} p{5.4cm}}
\toprule
Model & Key fields & Notable fields / relations \\
\midrule
\texttt{Team} &
\texttt{id} (PK) &
\texttt{name} (not unique --- names may repeat across challenges),
\texttt{challengeId?} (FK); has many \texttt{TeamMember}. \\
\addlinespace
\texttt{UserChallenge} &
\texttt{id} (PK); \texttt{@@unique([userId, challengeId])} &
\texttt{userId} (FK), \texttt{challengeId} (FK) --- enrollment join table. \\
\addlinespace
\texttt{TeamMember} &
\texttt{id} (PK); \texttt{@@unique([userId, teamId])} &
\texttt{userId} (FK), \texttt{teamId} (FK) --- membership join table. \\
\addlinespace
\texttt{TeamRequest} &
\texttt{id} (PK); \texttt{@@index([challengeId, status])} &
\texttt{name}, \texttt{message?}, \texttt{status}
(\texttt{TeamRequestStatus}, default \texttt{PENDING}), \texttt{challengeId}
(FK), \texttt{requestedById} (FK), \texttt{reviewedById?} (FK),
\texttt{reviewedAt?}. \\
\bottomrule
\end{tabular}
\end{table}

\noindent\textbf{Enumerations.}
\texttt{ActionCategory} = \{\,TRANSPORT, WATER, ENERGY, RECYCLING, FOOD\,\};\
\texttt{TeamRequestStatus} = \{\,PENDING, APPROVED, REJECTED\,\}. The
\texttt{[challengeId, status]} index on \texttt{TeamRequest} accelerates the
admin queue, which filters pending requests by challenge.

% ─────────────────────────────────────────────────────────────────────────────
\section{Algorithm Explanations}
% ─────────────────────────────────────────────────────────────────────────────

\subsection{Dense Rank with Tie Handling (Leaderboards)}

Items are sorted by value (descending), breaking ties by name. Ranks increment
only when the value changes, so tied entries share a rank (1, 2, 2, 3).

\begin{verbatim}
sort items by value desc, then name asc
rank <- 0; prev <- null
for each item:
    if prev is null or item.value != prev:
        rank <- rank + 1
        prev <- item.value
    item.rank <- rank
\end{verbatim}

\subsection{Streak Calculation (on each action)}

\begin{verbatim}
todayStart     <- midnight(now)
yesterdayStart <- todayStart - 1 day
newStreak <- 1
if user.lastActionDate exists:
    last <- midnight(user.lastActionDate)
    if   last == todayStart:     newStreak <- user.streak       # already today
    elif last == yesterdayStart: newStreak <- user.streak + 1   # consecutive
    else:                        newStreak <- 1                 # gap -> reset
update user.streak <- newStreak, lastActionDate <- now
\end{verbatim}

\subsection{Duplicate-Submission Guard (10\,s window + serializable txn)}

To prevent double-logging from rapid taps or retries, creation runs inside a
\texttt{SERIALIZABLE} transaction that first checks for an identical action
within the last ten seconds.

\begin{verbatim}
cutoff <- now - 10 seconds
BEGIN TRANSACTION (SERIALIZABLE):
    existing <- Action where userId, challengeId, category, actionType,
                points, note, imageUrl all match AND createdAt >= cutoff
    if existing: throw DUPLICATE          # -> HTTP 409
    else:        create Action
COMMIT
# a Prisma P2034 serialization failure is also returned as HTTP 409
\end{verbatim}

\subsection{Unique Challenge-Code Generation}

\begin{verbatim}
do:
    code <- random base36 string, 6 chars, upper-cased
while a Challenge with that code already exists
\end{verbatim}

% ─────────────────────────────────────────────────────────────────────────────
\section{Deployment Plan}
% ─────────────────────────────────────────────────────────────────────────────

\begin{itemize}
    \item \textbf{Hosting / CI-CD.} The app is hosted on \textbf{Vercel}.
    Building and deploying are triggered by pushing to \texttt{main} (or via
    \texttt{vercel --prod}); Vercel installs dependencies, runs
    \texttt{next build}, and serves the result.
    \item \textbf{Database.} \textbf{Neon} serverless PostgreSQL. The schema is
    applied with \texttt{prisma db push}; the project intentionally has no
    migrations folder for this phase.
    \item \textbf{Authentication.} \textbf{Clerk}. Administrative privileges
    are granted by setting a Clerk user's \texttt{publicMetadata.role} to
    \texttt{"admin"}.
    \item \textbf{Required environment variables.}
    \texttt{NEXT\_PUBLIC\_CLERK\_PUBLISHABLE\_KEY},
    \texttt{CLERK\_SECRET\_KEY}, and \texttt{DATABASE\_URL} (plus the Clerk
    sign-in/up redirect URLs).
\end{itemize}

% ─────────────────────────────────────────────────────────────────────────────
\begin{thebibliography}{9}
% ─────────────────────────────────────────────────────────────────────────────

\bibitem{nextjs}
Vercel, ``Next.js Documentation,'' 2024. [Online]. Available:
\url{https://nextjs.org/docs}

\bibitem{clerk}
Clerk, ``Clerk Documentation,'' 2024. [Online]. Available:
\url{https://clerk.com/docs}

\bibitem{prisma}
Prisma, ``Prisma ORM Documentation,'' 2024. [Online]. Available:
\url{https://www.prisma.io/docs}

\bibitem{neon}
Neon, ``Neon Serverless Postgres Documentation,'' 2024. [Online]. Available:
\url{https://neon.tech/docs}

\bibitem{c4}
S. Brown, ``The C4 model for visualising software architecture,'' 2024.
[Online]. Available: \url{https://c4model.com/}

\bibitem{cisc489}
E. Adekoya, K. Alkhaleefa, and J. Dunlap, ``Employee Sustainability Challenge
App: Final Project Documentation,'' CISC489 Topics in Human-Computer
Interaction, University of St.\ Thomas, Dec.\ 2025.

\end{thebibliography}

\end{document}
